\section{Program structure}

The programs are written in Python 3.13 with the standard library only: every method is
implemented from scratch, without NumPy. Every task has its own folder \texttt{task1},
\ldots, \texttt{task5} with the same files:

\begin{itemize}
    \item the module of the method (\texttt{lu.py}, \texttt{tridiagonal.py}, \texttt{iterative.py},
    \texttt{rotation.py}, \texttt{qr.py});
    \item \texttt{main.py}, which reads \texttt{input.txt} next to it (or the file given as an
    argument), solves the task and prints the result: \texttt{python3 -m lab1.task3.main};
    \item \texttt{input.txt} with the data of the variant and \texttt{output.txt} with the expected
    output, which the tests compare with the actual one.
\end{itemize}

The file \texttt{lab1/main.py} runs all five tasks one after another. All tests are in the folder
\texttt{lab1/tests}: one file per method, the tests of \texttt{common.py} and the tests that run
every program and compare its output with \texttt{output.txt}. The shared module
\texttt{common.py} contains the matrix operations (a matrix is a list of rows, a vector is a list
of numbers), the class \texttt{InputReader} that reads the size, the numbers and the precision
from the input file, the output formatting with four digits after the decimal point and the
function \texttt{run}. It computes the whole result before printing it, so an error never leaves a
partial output; errors are printed as \texttt{error: <message>} with the exit code 1. All sums are
accumulated in explicit loops: the built-in \texttt{sum} uses compensated summation since
Python 3.12 and would change the last bits of the results.

\lstinputlisting[caption={main.py}]{../main.py}
\lstinputlisting[caption={common.py}]{../common.py}

\pagebreak
